Ce dossier retrace les principaux projets qui ont marqué mes deux années de mastère, en particulier mon alternance au SHOM. Ils m'ont confronté à des contextes très différents, allant de l'évolution d'une chaîne de production historique à la conception d'applications, au traitement de données, au DevOps et au pilotage d'une équipe projet.

Au-delà des technologies utilisées, ces expériences m'ont surtout appris à mieux analyser un besoin et un existant, à justifier mes choix, à tenir compte des contraintes réelles d'un projet et à prendre du recul sur les solutions mises en œuvre. Elles m'ont également montré l'importance de la qualité, de la documentation, des tests et de la transmission pour qu'un développement reste réellement exploitable dans le temps.

À l'issue de ce parcours, je souhaite poursuivre dans le développement applicatif, avec un intérêt particulier pour le back-end, l'architecture logicielle et les problématiques liées aux données. Je souhaite continuer à évoluer sur des projets où le développement ne consiste pas seulement à produire du code, mais aussi à comprendre un système, le fiabiliser et le faire évoluer de manière durable.
